% For each dataset, we provide a standardized evaluation protocol based on realistic use cases, proposing improvements to current evaluation metrics by sampling more negative edges based on temporal patterns of the graph and proper evaluation metrics\todo{ER: Repetition}

%scalability and performance\todo{ER: Not clear what performance means here} perspectives.
% For the dynamic link property prediction task, we propose improvements to current evaluation protocol and provide standardized evaluation protocols based on realistic use-cases. 

\begin{abstract}
We present the \emph{Temporal Graph Benchmark~(\name)}, a collection of challenging and diverse benchmark datasets for realistic, reproducible, and robust evaluation of machine learning models on temporal graphs. \name datasets are of large scale, spanning years in duration, incorporate both node and edge-level prediction tasks and cover a diverse set of domains including social, trade, transaction, and transportation networks. For both tasks, we design evaluation protocols based on realistic use-cases. We extensively benchmark each dataset and find that the performance of common models can vary drastically across datasets. In addition, on dynamic node property prediction tasks, we show that simple methods often achieve superior performance compared to existing temporal graph models. We believe that these findings open up opportunities for future research on temporal graphs. Finally, \name provides an automated machine learning pipeline for reproducible and accessible temporal graph research, including data loading, experiment setup and performance evaluation. \name will be maintained and updated on a regular basis and welcomes community feedback. \name datasets, data loaders, example codes, evaluation setup, and leaderboards are publicly available at \url{https://tgb.complexdatalab.com/}.
\end{abstract}

%We present the \emph{Temporal Graph Benchmark~(\name)}, a collection of challenging and diverse benchmark datasets for realistic, reproducible, and robust evaluation of machine learning models on temporal graphs. \name datasets are of large scale, spanning years in duration, incorporate both node and edge-level tasks and cover a diverse set of domains including social, trade, transaction, and transportation networks. For both tasks, we provide standardized evaluation protocols based on realistic use cases and propose improvements to current evaluation methods. We extensively benchmark each dataset and 
% and show that simple methods often achieve competitive performance compared to existing temporal graph models. The large scale nature of \name  also opens up opportunities for future research. Finally, \name provides an automated machine learning pipeline for reproducible and accessible temporal graph research, including data loading, experiment design, and performance evaluation. \name will be maintained and updated on a regular basis and welcomes community feedback. \name datasets, data loaders, example codes, evaluation setup, and leaderboards are publicly available at \url{http://tgb.mila.quebec/}.



%addressing limitations of current evaluation techniques including small number of negative samples and inadequate performance metrics

% We present the Temporal Graph Benchmark~(\name): a collection of challenging and diverse benchmark datasets for realistic, reproducible and robust evaluation for temporal graph research.\name is motivated by three major limitations of current evaluation setups: 1). \emph{limited graph sizes}: almost all benchmark datasets are small with less than 2 million edges, 2). \emph{too few negative samples}: current evaluation for link prediction only samples one negative edge per positive edge, this is highly unrealistic and 3). \emph{lacking features}: all benchmark datasets contains no edge features outside of a few with edge weights. In this work, we addresses these limitations by adding six new datasets ranging from 4 million to 67 million, orders of magnitude larger than existing ones. Then we improves the evaluation for link prediction by sampling hundreds of negative edges per single positive edge. Lastly, we include additional edge and node features where possible to enrich the datasets. 

